\chapter{Usar la red: imágenes, ONNX y ROS 2}
\label{cap:integrar}

\begin{enpocas}
Detectar en una foto, exportar la red a ONNX para la computadora a bordo y
conectarla a la simulación para que el Tarot 650 busque el incendio con
FuegoNet. Archivos: \cod{red_fuego/inferir.py}, \cod{red_fuego/exportar.py},
\cod{detector_yolo.py} y \cod{launch/simulacion.launch.py}.
\end{enpocas}

\section{Detectar en una imagen}

\paso{25} \textbf{Probar el detector.}
La clase \cod{Detector} carga \cod{mejor.pt}, aplica letterbox, corre la red,
aplica el umbral de confianza (no hay NMS) y regresa las cajas a las coordenadas
de la imagen original.

Al cargar los pesos, \cod{fusionar} pliega cada BatchNorm en su convolución.
Al inferir, BatchNorm es solo una transformación lineal fija por canal:
\begin{equation}
y = \gamma\,\frac{W x - \mu}{\sqrt{\sigma^2+\epsilon}} + \beta
  = \underbrace{\frac{\gamma}{\sqrt{\sigma^2+\epsilon}}W}_{W'}\,x
  + \underbrace{\beta - \frac{\gamma\,\mu}{\sqrt{\sigma^2+\epsilon}}}_{b'},
\end{equation}
así que se puede meter dentro de los pesos de la convolución. Da exactamente las
mismas cajas (el mAP no cambia) con una operación menos por capa.

\begin{lstlisting}[style=bash]
~/.venvs/yolo/bin/python -m cuadrirotor_sim.red_fuego.inferir \
  ~/datasets/fuegonet/n_150/mejor.pt ~/datasets/incendio_bosque/images/val/00420.jpg salida.jpg
\end{lstlisting}

\begin{resultado}
\begin{lstlisting}
2 detecciones en 6.1 ms
  fuego  0.96  caja 117.2 322.9 255.2 476.0
  humo   0.96  caja 0.0 94.8 308.8 401.7
\end{lstlisting}
Cada fila es una caja: clase, confianza y $(x_1, y_1, x_2, y_2)$ en píxeles de la imagen
original de $640\times480$.
\end{resultado}


\section{Exportar a ONNX}

ONNX es un formato abierto que entienden TensorRT (Jetson), OpenVINO (Intel),
ONNX Runtime (cualquier CPU) y muchos otros. Como FuegoNet no tiene NMS, todo el
detector cabe en el grafo y la salida es un solo tensor $(1, 300, 6)$.

\paso{26} \textbf{Instalar \cod{onnx} sin romper NumPy.}
\begin{lstlisting}[style=bash]
~/.venvs/yolo/bin/pip install onnx onnxruntime
~/.venvs/yolo/bin/pip install "numpy==1.26.4" "ml-dtypes<0.6"   # ¡regresar NumPy!
\end{lstlisting}

\begin{leccion}
Al instalar \cod{onnx}, pip subió NumPy de 1.26 a 2.5 y \cod{ml-dtypes} a una
versión que exige NumPy~2. ROS 2 Jazzy está compilado contra NumPy~1 y el nodo
detector habría dejado de arrancar. La segunda línea regresa NumPy y fija una
\cod{ml-dtypes} compatible. Comprueba siempre con
\cod{python -c "import numpy; print(numpy.__version__)"} después de instalar algo.
\end{leccion}

\paso{27} \textbf{Exportar y comprobar.}
\begin{lstlisting}[style=bash]
~/.venvs/yolo/bin/python -m cuadrirotor_sim.red_fuego.exportar \
  ~/datasets/fuegonet/n_150/mejor.pt ~/datasets/fuegonet/fuegonet_n.onnx \
  ~/datasets/incendio_bosque/images/val/00420.jpg
\end{lstlisting}
El tercer argumento es una foto para comprobar que ONNX Runtime y PyTorch dan las
mismas detecciones.
\begin{resultado}
\begin{lstlisting}
Exportado /home/jairo/datasets/fuegonet/fuegonet_n.onnx
2 detecciones con confianza >= 0.25; diferencia máxima ONNX vs PyTorch: 3.05e-05
\end{lstlisting}
El archivo pesa 9.6~MB. La entrada se llama \cod{imagen} $(1,3,640,640)$ y la
salida \cod{detecciones} $(1,300,6)$.
\end{resultado}


\begin{leccion}
La primera comprobación usaba una entrada de ruido y comparaba las 300 filas. Daba
diferencias de más de 500 píxeles. No era un error de exportación: con ruido casi
todas las confianzas son $\approx 0$ y están empatadas, y \cod{topk} las ordena
distinto en PyTorch y en ONNX Runtime, así que se comparaban cajas de celdas
distintas. Ahora se usa una foto real y solo las detecciones con confianza
$\ge 0.25$. La diferencia baja a $3\times10^{-5}$ píxeles.
\end{leccion}

\section{Conectar FuegoNet a la simulación}

El nodo \cod{detector_yolo} ya publicaba \cod{/cuadrirotor/detecciones}
(\cod{vision_msgs/Detection2DArray}) con YOLO26 de Ultralytics. Para no duplicar
el nodo, \cod{inferir.py} incluye una envoltura \cod{ComoYolo} con la misma
interfaz que usa el nodo (\cod{.names} y \cod{.predict(...)[0].boxes}). El nodo
elige según la terminación del archivo de pesos:

\begin{lstlisting}[style=py]
        if v['modelo'].endswith('.fuegonet.pt'):
            from cuadrirotor_sim.red_fuego.inferir import ComoYolo
            self.modelo, self.etiqueta = ComoYolo(v['modelo'], v['dispositivo']), 'FuegoNet'
        else:
            self.modelo, self.etiqueta = YOLO(v['modelo']), 'YOLO26'
\end{lstlisting}

\paso{28} \textbf{Instalar los pesos en el paquete.}
\begin{lstlisting}[style=bash]
cp ~/datasets/fuegonet/n_150/mejor.pt \
   ~/ros2_ws/src/cuadrirotor_sim/yolo/incendio_fuegonet_n.fuegonet.pt
cd ~/ros2_ws && colcon build --packages-select cuadrirotor_sim --symlink-install
source install/setup.bash
\end{lstlisting}

\paso{29} \textbf{Volar la búsqueda con FuegoNet.}
El \emph{launch} tiene un argumento nuevo, \cod{red}, que elige el detector
(\cod{yolo26}, el de siempre, o \cod{fuegonet}):
\begin{lstlisting}[style=bash]
ros2 launch cuadrirotor_sim simulacion.launch.py mision:=busqueda red:=fuegonet
\end{lstlisting}
\begin{resultado}
En la terminal aparece \cod{YOLO listo: incendio_fuegonet_n.fuegonet.pt, clases
\{0: 'fuego', 1: 'humo'\}} y el panel «Detección YOLO» de Gazebo muestra
\cod{FuegoNet  xx ms} en la esquina. El resto de la misión (barrido, verificación,
aviso GPS) no cambia, porque los mensajes son los mismos.
\end{resultado}

\begin{consejo}
Para comparar las dos redes en las mismas condiciones, corre la misión dos veces
con el mismo \cod{mundo}, \cod{foco_x}, \cod{foco_y} y \cod{viento}, una con
\cod{red:=yolo26} y otra con \cod{red:=fuegonet}, y compara el tiempo hasta la
confirmación y las falsas alarmas.
\end{consejo}
